\chapter{Estado del arte y marco teórico}
\label{cap:estado-arte-marco-teorico}

\instruccion{Este capítulo complementa, pero no repite, los antecedentes del
problema presentados en la Introducción. Primero muestra qué investigaciones y
soluciones existen alrededor del problema; después desarrolla los fundamentos
con los que se comprenderá el objeto de estudio y se interpretarán los
resultados. Esta organización sigue siendo compatible con el desarrollo de la
perspectiva teórica de Hernández-Sampieri y Mendoza, pues integra la revisión
analítica de la literatura con la construcción del marco teórico
\cite{hernandezSampieri2018metodologia}.}

\section{Estado del arte y trabajos relacionados}

Analice investigaciones, métodos, sistemas o soluciones relacionados con la
tesis. Organice la literatura por enfoques, familias de soluciones, contextos,
variables o criterios de comparación; evite redactar una sucesión de resúmenes
independientes. Para los trabajos más relevantes considere propósito, contexto,
diseño metodológico, datos o participantes, propuesta, forma de evaluación,
resultados y limitaciones.

\instruccion{Al inicio de esta sección explique brevemente cómo localizó y
seleccionó la literatura: fuentes consultadas, periodo cubierto, términos o
cadenas relevantes, criterios de selección y fecha de la última búsqueda. Esto
no necesita convertirse en una subsección obligatoria. Si realizó una revisión
sistemática o un mapeo con un protocolo reproducible, presente aquí una síntesis
y traslade cadenas completas, diagramas de selección y listas extensas a un
apéndice. Para citar una fuente use
\texttt{\textbackslash cite\{kitchenham2007guidelines\}}, por ejemplo
\cite{kitchenham2007guidelines}.}

\instruccion{Compare los trabajos mediante criterios comunes para identificar
tendencias, coincidencias, contradicciones, limitaciones y aspectos todavía no
resueltos. Explique cómo esa evidencia se relaciona con el problema y con la
aportación de la tesis, sin crear una subsección separada sólo para repetir la
brecha de investigación.}

\ifmostrarEjemplos
\begin{table}[H]
\centering
\caption{Ejemplo de matriz comparativa de trabajos relacionados}
\label{tab:comparacion-trabajos}
\small
\begin{tabularx}{\textwidth}{L{2.4cm}YYYY}
\toprule
\textbf{Trabajo} & \textbf{Contexto y objetivo} & \textbf{Método o propuesta} & \textbf{Evaluación} & \textbf{Limitación relevante} \\
\midrule
Wohlin \cite{wohlin2014guidelines} & Describa el escenario. & Resuma el enfoque. & Indique datos y métricas. & Explique la limitación. \\
Autor B & Describa el escenario. & Resuma el enfoque. & Indique datos y métricas. & Explique la limitación. \\
\bottomrule
\end{tabularx}
\end{table}
\fi

\section{Fundamentos teóricos y conceptuales}

Desarrolle las teorías, modelos, enfoques, principios y conceptos necesarios
para comprender el estudio. Organice esta sección mediante subtítulos propios
del tema de tesis; no es obligatorio conservar una división genérica entre
teoría, conceptos, ingeniería de software o tecnología.

\instruccion{Los fundamentos teóricos explican relaciones, mecanismos o
principios mediante teorías, modelos y enfoques. Los fundamentos conceptuales
precisan el significado con el que se usarán términos especializados. En una
tesis de ingeniería de software ambos suelen estar estrechamente relacionados,
por lo que pueden desarrollarse juntos bajo los subtemas que el autor necesite.
Cuando existan definiciones rivales, compárelas y justifique la adoptada.}

\instruccion{Incluya fundamentos de ingeniería de software o tecnológicos sólo
cuando sean necesarios para comprender el estudio, sustentar una decisión,
definir una variable o establecer un criterio de evaluación. No convierta esta
sección en un manual de herramientas ni describa aquí versiones y
configuraciones específicas del proyecto; esos detalles corresponden a
Metodología.}

\instruccion{Cierre el capítulo mostrando para qué se utilizarán los fundamentos
seleccionados: formular o precisar la hipótesis, definir variables o categorías,
orientar el diseño de la solución, establecer criterios de evaluación o
interpretar los resultados. Cada fundamento incluido debe cumplir al menos una
de estas funciones.}

% Ejemplos de organización. Sustituya estos títulos por los temas reales:
% \subsection{Modelo o enfoque central del estudio}
% \subsection{Conceptos asociados al objeto de estudio}
% \subsection{Principios de ingeniería de software aplicables}
% \subsection{Fundamentos necesarios para la evaluación}
